\section{Higher-Order Language}
\label{sec:language}

To demonstrate semiring provenance tracking via automatic differentiation for higher-order programs, we define a simple total functional
programming language, extending the simply-typed lambda calculus. The language is parameterised by a signature
$\Sigma = (\PrimTy, \PrimOp)$ consisting of a set $\PrimTy$ of base types $\rho$ and a family of sets
$\PrimOp^\rho_{\rho_1,\ldots,\rho_n}$ of primitive operations $\phi$ of arity $n$ over those base types.

\subsection{Syntax}
\label{sec:language:syntax}

The syntax is defined in \figref{syntax}. Types includes base types $\rho$ drawn from $\PrimTy$, along with
standard type formers for sums, products, functions and lists% , plus a type constructor $\tyLift$ to allow
% approximation points to be added explicitly, as discussed in \secref{models-of-total-gps}
. Terms include
variables, the usual introduction and elimination forms% , a monadic return and bind for lifted terms
, and
primitive operations $\phi$.

The language is intentionally minimal: it excludes general recursion, and general inductive or coinductive
types, which we will consider in future work (\secref{conclusion}). Typing judgments for terms are standard
and shown in \figref{typing}, with the usual rules for products, sums, functions, and lists.

\input{fig/syntax}
\input{fig/typing}

\subsection{Semantics}
\label{sec:language:semantics}

\input{fig/semantics}

An interpretation of a signature $\Sigma = (\PrimTy, \PrimOp)$ can be given in any category $\cat{C}$ with
finite products, and assigns to each base type $\rho \in \PrimTy$ an object $\sem{\rho}_{\PrimTy}$ in
$\cat{C}$, and to each primitive operation $\phi \in \PrimOp^\rho_{\rho_1,\ldots,\rho_n}$, a morphism
$\sem{\phi}_{\Op}: \sem{\rho_1}_{\PrimTy} \times \ldots \times \sem{\rho_n}_{\PrimTy} \to
\sem{\rho}_{\PrimTy}$.

Assuming that $\cat{C}$ is bicartesian closed and has a list object (\eqnref{lists}), then we can extend an
interpretation of a signature $\Sigma$ to an interpretation of the whole language over
$\Sigma$. \figref{semantics:types} and \figref{semantics:contexts} define the interpretation of types and
contexts as objects of $\cat{C}$ respectively. Terms are interpreted as morphisms between the interpretations
of the context and type, as defined in \figref{semantics:terms}. We have used the notations $\pi_i$ for
projections, $\prodM{f}{g}$ for pairing, $\coprodM{f}{g}$ for (parameterised) copairing, $!_X$ for morphisms
to the terminal object, and $\lambda$ and $\eval$ for currying and evaluation for exponentials.

We have another interpretation $\sem{-}_{\mathit{fo}}$ of the
first-order types (those constructed from primitive types, sums, unit
and products) in any bicartesian category. Such interpretations are
preserved by finite coproduct and coproduct preserving functors:

\begin{lemma}\label{lem:first-order-agreement-types}
  If $\cat{C}$ and $\cat{D}$ are bicartesian and bicartesian closed categories with interpretations of the
  signature $\Sigma$, $F : \cat{C} \to \cat{D}$ is a bicartesian functor, and
  $F(\sem{\rho}_{\PrimTy}) \cong \sem{\rho}_{\PrimTy}$ for all $\rho$, then for all first-order types $\tau$,
  $F(\sem{\tau}_{\mathit{fo}}) \cong \sem{\tau}$, and similar for contexts only containing first-order types.
\end{lemma}

\subsubsection{Interpretation in the Semimodule Model}
\label{sec:language:gps-interpretation}

Given the above, we can now interpret the language in any of the bicartesian closed categories with list
objects we constructed in \secref{models-of-total-gps}. Specifically, we assume that we have an interpretation
of our chosen signature in $\Fam(\SDSemiMod(S))$. Each base type is
assigned a self-dual approximation object, and the first-order type formers preserve self-duality, so every
first-order type denotes a self-dual semimodule. Signatures are
first-order, so it does not matter that $\Fam(\SDSemiMod(S))$ is not closed. Any
such interpretation can be transported to $\Fam(\SemiMod(S))$ along the functor $H$ from \propref{ho-embedding} because it preserves finite products. We
can then interpret the whole language in $\Fam(\SemiMod(S))$.

Interpreting a program $\Gamma \vdash t : \tau$ yields a morphism of $\Fam(\SemiMod(S))$, the
underlying function together with its forward derivative. In contrast to many other bidirectional program
analysis techniques, a single morphism is enough, because a linear map has a transpose for free. At first-order types, \lemref{first-order-agreement-types} identifies the
interpretation with that of the first-order model $\Fam(\SDSemiMod(S))$, whose objects are self-dual, so the
forward derivative is a morphism of $\Fam(\SDSemiMod(S))$ whose conjugate $\conj{f}$ is the backward map
between the same semimodules. 

% Given a model of $\Sigma$ in $\Fam(\LatGal)$, the category $\Fam(\MeetSLat \times \JoinSLat^\op)$ can
% implement $\Sem$. It is bicartesian closed by \corref{mslat-jslat-bcc}, and can model $\tyList$ using infinite
% coproducts (\secref{models-of-total-gps:coproducts-and-products}). The embedding $\HoEmbed: \Fam(\LatGal) \to
% \Fam(\MeetSLat \times \JoinSLat^\op)$ from \propref{ho-embedding} preserves finite products, and so can
% transport each $\sem{\rho}_{\PrimTy}$ and $\sem{\phi}_{\Op}$ into the higher-order setting to provide a model
% of $\Sigma$. Finally the lifting monad in $\Fam(\LatGal)$ (\secref{first-order:lifting}) is also preserved by
% $H$. \roly{don't think we actually show this}

% The correctness of this interpretation is considered in the next section.

